% Revised 28 September 2026 using the corrected 24 September snapshot.
% Read METHODS_NOTES.md before submission. In particular, training and inference
% metadata and the provenance of the larger reference cohort remain incomplete.
\section{Methods}
\label{sec:methods}

\subsection{Benchmark datasets and molecular identity}

We evaluated conformer ensembles against experimental ligand geometries using
two complementary settings: recovery of one bound conformation per
protein--ligand complex and coverage of several observed conformations of the
same molecule. The CASF-2016 core and larger reference ligand collections
are denoted core and ref, respectively.\cite{su2019casf}
% AUTHOR: establish the exact source release and construction of CASF16_REF;
% the CASF citation alone does not document that local collection.
Ligands were matched to ChEMBL3D by exact canonical, heavy-atom isomeric SMILES.
ChEMBL3D provided molecular topologies and a separate computed conformer
ensemble for comparison.\cite{nikitin2025loqi}
Entries were retained when the ligand could be parsed, a matching topology was
available, and at least one eligible torsion was identified by the SMARTS
pattern \texttt{[!\$(*\#*)\&!D1]-!@[!\$(*\#*)\&!D1]}.
This torsion eligibility rule differs from the descriptor for the number of rotatable bonds used to stratify the results.

We resolved stereoisomers within each ChEMBL3D parent record using isomeric
SMILES and the tetrahedral and double-bond stereochemistry layers of InChI.
The selected topology record was retained in the mapping, and stored
coordinates were filtered to the matching stereoisomer. Thus, both the
reference ensemble and its conformer count were specific to the mapped
stereoisomer. The resulting core cohort contained 94 complex entries and 94
ChEMBL parent compounds; ref contained 1,236 entries and 1,042 parent compounds.
The cohorts shared no complex identifiers but included 17 common parent
compounds. We retained separate complex entries because they provide separate
experimental reference geometries, while accounting for repeated compounds
in the uncertainty analysis.

The second setting used a curated set of 23 druglike molecules with 2,450
stored experimental reference conformations in total, ranging from 23 to 716
per molecule. We used the supplied molecular identities and reference
coordinates without treating the frequency of a structure in this collection
as an equilibrium population. The collection is an incomplete set of observed conformations.
% AUTHOR: add the dataset release, selection and structure-cleaning procedure,
% and original dataset citation after verifying the supplied collection's lineage.

\subsection{Conformer generators}

The main comparison included four classical generators, five external learned
generators, and one representative Qwen checkpoint. The classical
generators comprised RDKit distance geometry and torsion perturbation, each
with and without energy minimization. RDKit conformers were generated with
ETKDGv3 from a molecular topology without coordinates, using random starting coordinates,
enforced chirality, and no RMSD pruning.\cite{wang2020etkdg}
The minimized variant applied MMFF94s for up to 500 iterations. Structures
that reached the iteration limit were retained if the optimization returned
normally; convergence was not required. The two variants used independently
sampled pools.

Torsion perturbation started from the matched ChEMBL3D topology's stored
geometry. In each trial, every eligible torsion received an independent
uniform angular displacement between $-120^{\circ}$ and $120^{\circ}$.
Candidates were rejected if a nonbonded heavy-atom distance was below 0.7
times its distance-geometry lower bound; bonded and 1,3 pairs were excluded
from this check. The minimized variant additionally applied MMFF94s and
repeated the clash check. Candidate generation continued toward the requested
pool size before the common validity filter was applied.

The external methods were LoQI,\cite{nikitin2025loqi} Torsional
Diffusion,\cite{jing2022torsional} Molecular Conformer Fields (MCF,
drugs-L),\cite{wang2024mcf} NExT-Mol
(DMT-L),\cite{liu2025nextmol} and FLOWR.root, referred to here as
FlowR.\cite{cremer2025flowr} We evaluated their saved conformer outputs through
the benchmark adapters. FlowR used the v2.2 checkpoint for generation of ligands alone, with the
input atom and bond graph fixed during coordinate generation. Its adapter
included up to 200 steps of force-field cleanup, using MMFF with UFF as a
fallback, and stereochemistry checks; mirror reflection was allowed when it
restored the requested stereoisomer. Remaining stereochemical mismatches
were rejected. Consequently, an ensemble labeled raw in the benchmark can
already include preprocessing specific to its generator.

The Qwen generator produced conformations autoregressively from molecular
SMILES. The main analysis used the model with 1.7 billion parameters and finite
scalar quantization (FSQ), evaluated at pretraining step 47,023. The
comparison with multiple experimental references used the same representative Qwen checkpoint and the five external methods.
Overlap of the training set with the benchmark molecules has not been established
by a common audit.
% AUTHOR: document the checkpoint-selection rule, including any use of these
% benchmark outcomes. Supply training datasets/splits, representation and decoder
% details, inference settings, checkpoint hashes, and Qwen/FSQ citations.

\subsection{Ensemble sizes and validity filtering}

The main CASF comparison used two candidate targets per ligand, both defined
before the common validity filter. The ChEMBL-count tier used a target equal
to the stored ChEMBL3D conformer count for the mapped stereoisomer, while the
fixed tier used a target of 1,000 conformers. The ChEMBL-count tier was sampled
without replacement from each available candidate pool. Subsets were capped
at the available pool size and then subjected to the same validity filter.
The subsets were drawn from existing outputs, without separate generation
runs or matching of computational cost.
The launcher for classical generation used base seed 42; the common
procedure for assembling ensembles from saved pools used base seed 1729.
Seeds for each molecule and method were derived from these values.

We applied PoseBusters to assess molecular identity and intramolecular
plausibility.\cite{buttenschoen2024posebusters} The configured checks covered
loading, sanitization, InChI conversion, connectivity, radicals, molecular
formula, bond connectivity, tetrahedral and double-bond stereochemistry,
bond lengths and angles, internal clashes, planarity, and energy ratio.
The distance-geometry tolerances were 0.25 for bond lengths and angles and
0.3 for clashes, with hydrogens excluded from these distance checks.
Planarity tolerance was 0.25~\AA{} for the specified aromatic rings and
trigonal carbon--carbon double-bond groups. The threshold for the energy ratio was
100, using a reference ensemble of 50 conformations. No checks of contacts with the protein were included.

A conformer passed only if all configured Boolean checks passed. Generated
CASF ensembles were checked against the matched ChEMBL3D molecular reference;
the stored ChEMBL3D comparison ensembles were checked against the CASF ligand.
The same pass mask was applied to the full candidate pool and its sampled
subsets. Rejected conformers were not replaced after filtering. Retained
ensemble sizes and numbers of generation attempts could differ between
methods at the same candidate target.
ChEMBL3D was evaluated as a finite stored ensemble, with a cap of 2,000 conformers for the designated large reference entry, 1tlp, when necessary.

\subsection{Conformational recovery and diversity}

CASF recovery was evaluated on conformers that passed PoseBusters, using the
experimental ligand geometry as the target. We removed hydrogens and used
RDKit's symmetry-aware \texttt{GetBestRMS} alignment. If this alignment raised
an exception, the CASF implementation attempted \texttt{AlignMol} with atoms matched by their indices. An unresolved, nonfinite RMSD caused analysis of
that ensemble to fail.

For each ligand, we calculated the minimum and median RMSD over the retained
ensemble. Hit@$\delta$ was one when at least one conformer had RMSD
$\leq\delta$, and zero otherwise, at thresholds of 0.25, 0.5, 0.75, and
2.0~\AA. We used 0.75~\AA{} as the principal recovery cutoff
and the remaining thresholds to assess sensitivity. Cohort hit rates used all mapped entries as the denominator,
including empty or missing outputs as misses. We averaged the minimum RMSD, ensemble median RMSD, and diversity
measures only over entries with defined values.

We measured geometric diversity by greedy clustering of the retained
conformers at a principal heavy-atom RMSD threshold of 1.0~\AA{},
with additional radii of 0.5 and 2.0~\AA{} for sensitivity analysis. Conformers were visited
in stored order and assigned to the nearest existing cluster representative
if their RMSD was below the threshold; otherwise, a new cluster was created.
We reported the mean number of clusters across measurable ensembles.
This measure describes geometric spread and depends on ensemble size
and clustering order.

\subsection{Coverage of multiple experimental conformations}

For each druglike molecule, we compared the supplied generated pool with all
of its experimental references using aligned heavy-atom RMSD. Recall
coverage (COV-R) was the fraction of reference conformations with a generated
neighbor below 0.75~\AA. Precision coverage (COV-P) was the fraction of
generated conformations with a reference neighbor below the same threshold.
Matching distance MAT-R was the mean RMSD from each reference to its
nearest generated conformer; MAT-P was the mean RMSD from each generated
conformer to its nearest reference. Unlike CASF Hit@$\delta$, coverage used a strict inequality.
We averaged each metric equally across molecules. COV-P measures proximity
to the available observations; a low value does not establish that the
remaining conformers are physically inaccessible.

These coverage calculations used the supplied generated pools without a
common PoseBusters filter. Validity was assessed separately. The imported
Qwen evaluator counted reference or generated conformations with no finite
RMSD match as coverage misses, whereas the local evaluator for external methods
excluded them from the corresponding coverage denominator. Both omitted
undefined distances to the nearest match from MAT. Differences in filtering
and failure handling limit the comparability of these coverage estimates.

\subsection{Statistical analysis}

We calculated paired differences in recovery and coverage on the same
complex entries or druglike molecules. For CASF, 95\% confidence intervals were obtained from
3,000 bootstrap resamples of ChEMBL parent compounds, retaining all complex
entries belonging to each sampled compound. Each resample used the mean
over its complex entries, preserving equal weighting of entries while
accounting for repeated compounds. Missing hits were assigned zero.
For druglike method differences, we used 10,000 paired resamples of the
23 molecules. Intervals for each method's mean COV-R and COV-P
(Figure~\ref{fig:druglike}A) used 10,000 resamples of the same molecules.

Intervals were the 2.5th and 97.5th percentiles of the bootstrap distribution,
using a NumPy random generator initialized with seed 20260924. These
intervals quantify uncertainty from resampling the evaluated molecules.
They do not include variation across independent generation or training runs.
% Provenance: CASF/drug paired intervals were calculated in
% docs/publication_analysis_2026_09_24.py and archived in paired_cluster_bootstrap.csv
% and druglike_paired.csv. Figure 4 marginal intervals were added by
% manuscript/build_figures.py, not provided by the public dashboard; see
% figure-data/drug_coverage_intervals.csv.

As an overlap sensitivity analysis, we repeated ref comparisons after
excluding the 17 parent compounds also present in core, leaving 1,189
complex entries from 1,025 compounds. Exploratory size analyses used
counts of heavy atoms from the experimental ligand after removing hydrogens, grouped
as fewer than 20, 20--29, 30--39, and at least 40 atoms. Flexibility analyses
used RDKit counts of rotatable bonds from that same reference, grouped as 0--3,
4--6, 7--8, and at least 9. These descriptors from the experimental reference assigned each
entry to the same group for every generator. Stratum hit rates included all
mapped entries in the group, with missing outputs counted as misses.
These descriptive analyses did not adjust size effects for flexibility or
chemical class. Confidence intervals were not adjusted for multiple
comparisons or checkpoint selection.
